$!A \ident (!B \lif !C)$ হলে $\pi$-র শেষাংশ
\[
\AxiomC{}
\RightLabel{$\pi_1'$}
\Deduce$!B, \Gamma \fCenter \Delta, !C$
\RightLabel{\RightR{\lif}}
\UnaryInf$\Gamma \fCenter \Delta, !B \lif !C$
\LeftSubproofLabel{$\pi_1$}
\AxiomC{}
\RightLabel{$\pi_2'$}
\Deduce$!B \lif !C, \Pi \fCenter \Lambda, !B$
\AxiomC{}
\RightLabel{$\pi_2''$}
\Deduce$!C, \Pi \fCenter \Lambda$
\RightLabel{\LeftR{\lif}}
\BinaryInf$!B \lif !C, \Pi \fCenter \Lambda$
\RightSubproofLabel{$\pi_2$}
\RightLabel{\Cut}
\BinaryInf$\Gamma, \Pi \fCenter \Delta, \Lambda$
\DisplayProof
\]
নীচের \Cut{}-এ শেষ হওয়া !!{proof}টি দেখি:
\[
\AxiomC{}
\RightLabel{$\pi_1'$}
\Deduce$!B, \Gamma \fCenter \Delta, !C$
\RightLabel{\RightR{\lif}}
\UnaryInf$\Gamma \fCenter \Delta, !B \lif !C$
\LeftSubproofLabel{$\pi_1$}
\AxiomC{}
\RightLabel{$\pi_2'$}
\Deduce$!B \lif !C, \Pi \fCenter \Lambda, !B$
\RightLabel{\Cut}
\BinaryInf$\Gamma, \Pi \fCenter \Delta, \Lambda, !B$
\DisplayProof
\]
এর কর্তন-!!{height} $\pi$-র চেয়ে কম। তাই আরোহের অনুমান
$\Gamma, \Pi \fCenter \Delta, \Lambda, !B$-র
!!a{proof}~$\pi_1''$ দেয়। এবার এই \Cut{}-এ শেষ হওয়া
!!{proof}টি দেখি:
\[
\AxiomC{}
\RightLabel{$\pi_1'$}
\Deduce$!B, \Gamma \fCenter \Delta, !C$
\AxiomC{}
\RightLabel{$\pi_2''$}
\Deduce$!C, \Pi \fCenter \Lambda$
\RightLabel{\Cut}
\BinaryInf$!B, \Gamma, \Pi \fCenter \Delta, \Lambda$
\DisplayProof
\]
% BN-SRC-531 TARGET-CORRECTION-BEGIN
\par\smallskip\noindent\textnormal{\small\textbf{উৎস-সংশোধন (BN-SRC-531):} এই প্রমাণবৃক্ষের বাম পূর্বধারণায় উল্টানো `B!`-এর বদলে আগের বিধির `!B` বসানো হয়েছে।}
% BN-SRC-531 TARGET-CORRECTION-END
এর কর্তনমাত্রা ও কর্তন-!!{height} দুটিই $\pi$-র সংশ্লিষ্ট
কর্তনমাত্রা ও !!{height}-র চেয়ে কম। তাই আরোহের অনুমানে এর অন্তিম
সিকোয়েন্টের একটি \Log{G3c}-প্রমাণ পাই; পরে দুর্বলীকরণের
গ্রহণযোগ্যতায় অতিরিক্ত প্রসঙ্গ যোগ করে
$!B, \Gamma, \Pi, \Pi \fCenter \Delta, \Lambda, \Lambda$-র
\Log{G3c}-!!{proof}~$\pi_2'''$ পাই।
% BN-SRC-532 TARGET-CORRECTION-BEGIN
\par\smallskip\noindent\textnormal{\small\textbf{উৎস-সংশোধন (BN-SRC-532):} আরোহের অনুমান সরাসরি দ্বিতীয় কর্তনের অন্তিম সিকোয়েন্টের প্রমাণ দেয়; এখানে দেখানো অতিরিক্ত প্রসঙ্গ পেতে দুর্বলীকরণের মধ্যবর্তী ধাপটি উল্লেখ করা হয়েছে।}
% BN-SRC-532 TARGET-CORRECTION-END
এবার নীচের \Cut{}-এ শেষ হওয়া !!{proof}টিতে আরোহের
অনুমান প্রয়োগ করি:
\[
\AxiomC{}
\RightLabel{$\pi_1''$}
\Deduce$\Gamma, \Pi \fCenter \Delta, \Lambda, !B$
\AxiomC{}
\RightLabel{$\pi_2'''$}
\Deduce$!B, \Gamma, \Pi, \Pi \fCenter \Delta, \Lambda,
\Lambda$
\RightLabel{\Cut}
\BinaryInf$\Gamma, \Gamma, \Pi, \Pi, \Pi \fCenter \Delta, \Delta, \Lambda, \Lambda, \Lambda$
\DisplayProof
\]
% BN-SRC-533 TARGET-CORRECTION-BEGIN
\par\smallskip\noindent\textnormal{\small\textbf{উৎস-সংশোধন (BN-SRC-533):} শেষ কর্তনের সিদ্ধান্তে দুই পূর্বধারণার অবশিষ্ট প্রসঙ্গগুলি মিলিয়ে লেখা হয়েছে; কর্তনসূত্র সিদ্ধান্তে রাখা হয়নি।}
% BN-SRC-533 TARGET-CORRECTION-END
(এর কর্তনমাত্রা $< \cutrank{\pi}$।) ফলে
$\Gamma, \Gamma, \Pi, \Pi, \Pi \fCenter \Delta, \Delta, \Lambda, \Lambda, \Lambda$-র
!!a{proof} পাই।
% BN-SRC-534 TARGET-CORRECTION-BEGIN
\par\smallskip\noindent\textnormal{\small\textbf{উৎস-সংশোধন (BN-SRC-534):} শেষ সিদ্ধান্তের সঙ্গে মিলিয়ে গদ্যের উত্তরপক্ষে দ্বিতীয় \texttt{\textbackslash Delta} ফিরিয়ে আনা হয়েছে।}
% BN-SRC-534 TARGET-CORRECTION-END
% BN-SRC-535 TARGET-CORRECTION-BEGIN
\par\smallskip\noindent\textnormal{\small\textbf{উৎস-সংশোধন (BN-SRC-535):} অনির্ধারিত কর্তনমাত্রা-নির্দেশের বদলে সংজ্ঞায়িত নির্দেশ বসানো হয়েছে।}
% BN-SRC-535 TARGET-CORRECTION-END
\olref[seq][inv]{prop:G3c-cont-adm} প্রয়োগ করে
$\Gamma, \Pi \fCenter \Delta, \Lambda$-র !!{proof}~$\pi'$
পাই।
